%% kapa: exam=1 round=34 year=2023 period=1 form=A subject=civil_law subject_ko=민법 seq=17 number=17
%% source: Q-Net 감정평가사 기출문제 (cst00309) artlSeq=5211789 file=raw/1008/5211789/2023년도 제34회 감정평가사 1차시험 1교시.pdf
%% conversion: kordoc 4.24.51 -> tools/gen1.py (auto); review=proofread-sonnet; flags=none
\begin{question}{17}
법률행위의 취소에 관한 설명으로 옳지 않은 것은? (다툼이 있으면 판례에 따름)
\begin{choices}
  \item 취소권의 단기제척기간은 취소할 수 있는 날로부터 3년이다.
  \item 취소권의 행사시 반드시 취소원인의 진술이 함께 행해져야 하는 것은 아니다.
  \item 취소할 수 있는 법률행위의 상대방이 그 행위로 취득한 특정의 권리를 양도한 경우, 양수인이 아닌 원래의 상대방에게 취소의 의사표시를 하여야 한다.
  \item 노무자의 노무가 일정 기간 제공된 후 행해진 고용계약의 취소에는 소급효가 인정되지 않는다.
  \item 매도인이 매매계약을 적법하게 해제한 후에도 매수인은 그 매매계약을 착오를 이유로 취소할 수 있다.
\end{choices}
\end{question}
\answer{1}
